\verseline{}{他含泪说完，便放开缰绳，让舰队前行，}
\verseline{}{终于驶抵优波亚人所建库迈的海岸。}
\verseline{}{他们掉转船首，朝向大海；铁锚以咬紧}
\verseline{}{海底的利齿稳住船只，弯曲的船尾}
\verseline{5}{沿岸排开。一群青年急切地跃上}
\verseline{}{赫斯珀里亚海岸；一些人从燧石的纹理中}
\verseline{}{寻找深藏的火种，一些人奔入野兽}
\verseline{}{栖居的密林，指点新发现的溪流。}
\verseline{}{虔敬的埃涅阿斯却走向高处的圣所——}
\verseline{10}{崇高的阿波罗在那里主掌一切——又远远寻找}
\verseline{}{可畏西比尔隐居的巨大洞穴；提洛的先知}
\verseline{}{赋予她宏阔的心神，为她揭开未来。}
\verseline{}{他们已走进三岔路女神的树林与金色殿宇。\footnote{库迈是优波亚岛哈尔基斯人的希腊殖民城。埃涅阿斯先到阿波罗神庙，继而前往西比尔洞窟，准备下行冥界。“三岔路女神”是赫卡忒的称号，兼有道路交汇、夜间与幽冥的意味；原文称号 Trivia 在本卷统译为“三岔路女神”（其他卷音译为“特里维亚”），与“赫卡忒”交替称呼同一女神。}}
\verseline{}{相传代达罗斯逃离弥诺斯的王国，}
\verseline{15}{敢凭迅疾的羽翼把自己托付给天空，}
\verseline{}{沿无人走过的航路飞渡到寒冷的大熊星座下，}
\verseline{}{最后轻盈地停在哈尔基斯人的高城。}
\verseline{}{他一回到这片土地，便先向你，福玻斯，献上}
\verseline{}{那双以羽翼为桨的飞具，又建起宏大的神殿。}
\verseline{20}{门上刻着安德罗革俄斯之死；继而是被迫}
\verseline{}{偿付刑罚的雅典人——何其可悲！——每年献出}
\verseline{}{七名儿女；抽签已毕，签瓮立在一旁。}
\verseline{}{对面，克诺索斯的土地从海上高高升起，与之相应：}
\verseline{}{这里是对公牛的残酷迷恋，帕西法厄暗中藏进}
\verseline{25}{木牛，与异类交合；还有那混血、双形的后代}
\verseline{}{弥诺陶洛斯，邪恶爱欲留下的见证；}
\verseline{}{这里又是那座艰辛的宫室与不可解脱的迷途。}
\verseline{}{可是代达罗斯怜悯王女深重的爱情，}
\verseline{}{亲手解开宫殿的诡局与曲折，}
\verseline{30}{以丝线引导盲目的脚步。伊卡洛斯啊，}
\verseline{}{若悲痛容许，你也会在这浩大的作品中占很大一部分。}
\verseline{}{他两次试图用黄金刻出你的坠落；}
\verseline{}{两次，父亲的双手都垂落下来。\footnote{神庙门上的浮雕叙述代达罗斯神话：安德罗革俄斯死后，雅典人被迫向克里特进献童男童女；帕西法厄借木牛与公牛交合，生下弥诺陶洛斯；代达罗斯营造迷宫，又以丝线帮助阿里阿德涅引忒修斯脱困。浮雕本应刻下伊卡洛斯之死，代达罗斯的手却两度因父爱而坠落，不能完成作品。}他们本会立即}
\verseline{}{用目光把一切看完，若非先行的阿卡忒斯}
\verseline{35}{已经回来；与他同来的还有福玻斯与三岔路女神的祭司，}
\verseline{}{格劳库斯之女得伊福柏，她对国王说道：\footnote{得伊福柏是格劳库斯之女，即库迈西比尔。原诗按职能交替称她为“西比尔” “祭司” “先知”和“处女”，这些称谓均指同一人物。她先催促献祭；埃涅阿斯求取神谕、进入冥界之前，都先举行祭仪。}}
\verseline{}{“此刻的情势不容你观看这些图景；}
\verseline{}{现在最好从未曾负轭的牛群中宰杀七头公牛犊，}
\verseline{}{再依礼选出同样数目的两岁母羊。”}
\verseline{40}{她这样告诫埃涅阿斯；他的人没有耽搁}
\verseline{}{神圣的祭礼，女祭司便召特乌克罗斯人进入高耸的神殿。}
\verseline{}{优波亚岩山的一整面被凿成洞穴，}
\verseline{}{一百条宽阔的通道通往那里，还有一百扇门；}
\verseline{}{同样多的声音从中涌出，传来西比尔的答语。}
\verseline{45}{他们来到门槛前，少女喊道：“求问命运的}
\verseline{}{时辰到了；神来了，看哪，神来了！”她还在门前}
\verseline{}{这样说着，面容忽然不复先前，肤色也已改变，}
\verseline{}{梳好的头发散开；她的胸膛喘息，}
\verseline{}{凶猛的心因狂乱而胀起；她显得更加高大，}
\verseline{50}{声音也不再属于凡人——神已逼近，}
\verseline{}{以神力的气息吹入她身中。“特洛伊的埃涅阿斯，”}
\verseline{}{她说道，“你还迟迟不肯许愿、祈祷吗？还迟疑吗？}
\verseline{}{在此之前，震怖的殿宇不会张开巨口。”}
\verseline{}{她说罢便沉默。冰冷的战栗穿过特乌克罗斯人}
\verseline{55}{坚硬的骨骼，国王从胸膛深处倾出祈祷：}
\verseline{}{“福玻斯啊，你一向怜悯特洛伊沉重的苦难，}
\verseline{}{是你引导帕里斯那达耳达诺斯人的手与箭矢，}
\verseline{}{射中埃阿科斯后裔的身体；也由你引领，我走遍}
\verseline{}{环抱广大土地的无数海域，深入遥远的}
\verseline{60}{马叙利亚部族和叙耳提斯环绕的田野，}
\verseline{}{如今终于抓住那一再逃离我们的意大利海岸；}
\verseline{}{愿特洛伊的厄运只追随我们到这里为止。}
\verseline{}{所有曾与伊利乌姆、与达耳达尼亚的}
\verseline{}{赫赫荣光为敌的男神女神，如今也该依神律}
\verseline{65}{宽恕帕加马的子民了。而你，至圣的女先知，}
\verseline{}{预知将来，我并不求命运未许的王国；}
\verseline{}{请容许特乌克罗斯人定居拉提乌姆，}
\verseline{}{让特洛伊流亡的众神与久受颠沛的神灵安居。}
\verseline{}{那时我将以整块大理石为福玻斯与三岔路女神}
\verseline{70}{建起圣殿，并设立以福玻斯命名的节日。}
\verseline{}{在我们的王国里，也有广大的内殿等待着你：}
\verseline{}{我将在那里安放你的神签，以及你向我族}
\verseline{}{启示的隐秘命数，并选出专职的男子奉祀你，}
\verseline{}{慈惠者啊。只求你别把神歌托付给树叶，}
\verseline{75}{任它们被疾风吹乱，在空中受风戏弄；}
\verseline{}{请你亲口唱出。”说到这里，他止住了声音。\footnote{西比尔把神谕写在树叶上，排列于洞中；门一开，风便可能吹乱文字，使求问者无法恢复次序。埃涅阿斯因此请求她亲口宣示，并承诺在罗马为阿波罗和特里维亚设立圣所、选任祭司、保存神谕。这一承诺又使人联想到罗马后来保存的《西比尔神谕集》；库迈传说和后世制度分属不同传统。}}
\verseline{}{女先知还未屈服于福玻斯，仍在洞中凶猛地}
\verseline{}{狂奔挣扎，想把强大的神从胸膛里甩脱；}
\verseline{}{神却越发磨折她发狂的口，驯服她}
\verseline{80}{狂野的心，以不断的重压塑造她。}
\verseline{}{殿宇的一百扇巨门已经自行开启，}
\verseline{}{把女先知的答语送入风中：}
\verseline{}{“啊，你终于摆脱大海的重重险难！}
\verseline{}{陆上的灾祸却更为沉重。达耳达诺斯人将进入}
\verseline{85}{拉维尼乌姆的国土；从心中放下这一忧虑吧；}
\verseline{}{但他们也会希望自己从未来过。战争，可怖的战争！}
\verseline{}{我看见台伯河翻涌着重重血沫。}
\verseline{}{你不会缺少西摩伊斯、克珊托斯或多利安人的营垒；}
\verseline{}{拉提乌姆已经生出另一个阿喀琉斯，}
\verseline{90}{他同样出自女神；那始终缠住特乌克罗斯人的朱诺}
\verseline{}{也绝不会远离。到那时，你身处窘境，低身求援，}
\verseline{}{意大利哪个民族、哪座城邑会逃过你的哀求！}
\verseline{}{如此大祸的起因，又将是特乌克罗斯人的一位}
\verseline{}{外邦妻子；又将是一场异族婚姻。}
\verseline{95}{你不要向灾难屈服，要以更大的胆魄迎它而去，}
\verseline{}{无论你的命运准许你向何处前行。最初的救援之路，}
\verseline{}{会从你最意想不到的一座希腊城市开启。”\footnote{“另一个阿喀琉斯”通常指图尔努斯；“外国新娘”是拉维尼娅，她将再次成为战争的婚姻根源；“希腊人的城市”则预示厄万德统治的帕兰忒乌姆会成为特洛伊人的盟友。阿波罗附在西比尔身上时被比作骑手：以缰绳和刺棒控制她，使她不得不服从神意。}}
\verseline{}{库迈的西比尔从内殿唱出这些话，}
\verseline{}{言辞可怖而曲折，吼声回荡洞中，}
\verseline{100}{把真实裹在幽暗里；她发狂时，阿波罗}
\verseline{}{这样抖动她的缰绳，又在她胸中旋动刺棒。}
\verseline{}{狂乱刚退，发狂的口刚刚平静，}
\verseline{}{英雄埃涅阿斯便开口：“没有哪一种苦难，}
\verseline{}{少女啊，会以陌生或意外的面貌在我眼前升起；}
\verseline{105}{一切我都预先想过，早已在心中一一经历。}
\verseline{}{我只有一个请求：既然这里据说是地下君王的门户，}
\verseline{}{又有阿刻戎回流而成的幽暗沼泽，}
\verseline{}{愿我能前去，亲眼看见敬爱的父亲，面对他的容颜；}
\verseline{}{请教我那条路，为我打开神圣的门。}
\verseline{110}{我曾穿过烈火和身后追来的千百支兵矛，}
\verseline{}{用这双肩把他救出，从敌阵中央夺回；}
\verseline{}{他陪伴我走完航程，同我承受所有海域、}
\verseline{}{一切大海与天空的威胁；尽管衰弱，}
\verseline{}{他承担的远远超过老年所得的力量与定数。}
\verseline{115}{也是他嘱托我来到你的门槛，低身向你求助，}
\verseline{}{一再恳求，把这使命交给我。慈惠者啊，}
\verseline{}{我求你怜悯这父亲和儿子；你无所不能，}
\verseline{}{赫卡忒绝非徒然让你主掌阿韦尔努斯的树林。}
\verseline{}{若俄耳甫斯能够召回妻子的亡魂，}
\verseline{120}{凭色雷斯竖琴和鸣响的琴弦，}
\verseline{}{若波吕克斯能以轮番赴死赎回兄弟，}
\verseline{}{一次次沿这条路往返；我又何必提伟大的忒修斯，}
\verseline{}{何必提阿尔喀得斯？我的血脉也出自至高的朱庇特。”}
\verseline{}{他以这些话祈求，双手紧握祭坛；}
\verseline{125}{女先知这才开口：“生自众神血脉的后裔，}
\verseline{}{安喀塞斯之子、特洛伊人啊，下到阿韦尔努斯并不艰难；}
\verseline{}{黑暗狄斯的门日夜敞开；}
\verseline{}{可要折回脚步，逃回上界的空气，}
\verseline{}{这才是真正的难关，真正的考验。\footnote{阿韦尔努斯是库迈附近的火山湖，此处兼指冥界入口。西比尔强调：下行容易，返回人间才艰难。}只有少数人——公正的}
\verseline{130}{朱庇特所钟爱的人，或以炽烈德行升入以太的人，}
\verseline{}{众神的子孙——能够做到。森林占据中途的一切，}
\verseline{}{科库托斯以幽黑的河湾环绕四周。}
\verseline{}{倘若你心中有这样深的爱，这样强的渴望，}
\verseline{}{要两次渡过斯提克斯的湖，两次看见}
\verseline{135}{幽黑的塔耳塔洛斯，并乐于纵容这疯狂的劳苦，}
\verseline{}{那就听清首先必须完成的事。一株浓荫的树上}
\verseline{}{藏着一根金枝，叶片和柔韧的枝条皆为黄金，}
\verseline{}{被奉为地下朱诺的圣物；整座树林遮掩着它，}
\verseline{}{幽暗谷地的阴影把它封闭。}
\verseline{140}{然而，谁若不先从树上折取金发般的枝叶，}
\verseline{}{就无权进入大地掩藏的深处。}
\verseline{}{美丽的普洛塞庇娜规定，要把它作为自己的}
\verseline{}{礼物献来。第一枝折下，另一根金枝便接着生出，}
\verseline{}{枝条仍以同样的金属抽出新叶。}
\verseline{145}{所以高高举目，细细寻找；找到以后，}
\verseline{}{便依礼用手折取。若命运召唤你，它会自愿而轻易地}
\verseline{}{随你而去；否则，任何力量都无法胜过它，}
\verseline{}{坚硬的刀刃也不能把它劈落。\footnote{金枝是献给普洛塞庇娜的神圣礼物，也是命运是否准许埃涅阿斯入冥的征验。若命运召唤，它会自行顺从；若非天命，任何力量都不能折下。金枝因此同时是祭品、命运征验和入冥凭证。}}
\verseline{}{此外，一位友人的遗体正躺在那里——}
\verseline{150}{唉，你还不知道——未葬的死者使整支舰队蒙受不洁；}
\verseline{}{你却还在求问神谕，悬心立在我的门前。}
\verseline{}{先把他送回应有的居所，葬入坟墓。}
\verseline{}{再牵来黑色牲畜，作为最初的祓罪之礼；}
\verseline{}{只有这样，你才会看见斯提克斯的树林与}
\verseline{155}{活人不能踏入的国度。”她说完，抿紧嘴唇，沉默下来。\footnote{舰队中尚有一具未葬的尸体，死亡造成的仪式性污染必须先被清除；在此之前，埃涅阿斯不能接近冥界。西比尔命人安葬死者、建墓并以黑色牲畜赎罪，使寻枝、葬礼和净罪成为同一套入冥准备。}}
\verseline{}{埃涅阿斯神色悲哀，垂下目光，}
\verseline{}{离开洞穴前行，心中反复推想}
\verseline{}{幽暗未明的结局。忠实的阿卡忒斯}
\verseline{}{伴他同行，怀着同样的忧虑踩下脚步。}
\verseline{160}{他们彼此交谈，提出种种猜测：}
\verseline{}{女先知所说的死去同伴究竟是谁，哪具遗体必须安葬。}
\verseline{}{来到干燥的海岸，他们便看见米塞努斯}
\verseline{}{横死在那里，遭遇了不应有的死亡——}
\verseline{}{埃俄卢斯之子米塞努斯；没有别人更能用铜号}
\verseline{165}{唤起战士，以号音点燃战斗。}
\verseline{}{他曾是伟大赫克托耳的同伴，手执号角与长矛，}
\verseline{}{在赫克托耳身旁卓然出入一场场战斗；}
\verseline{}{胜利者阿喀琉斯夺走那人的生命以后，}
\verseline{}{这位勇武的英雄便加入达耳达诺斯人埃涅阿斯的行列，}
\verseline{170}{成为他的同伴——这番追随也不逊于从前。}
\verseline{}{可就在这时，他疯狂地吹响空心的海螺，}
\verseline{}{让号声传遍海面，以乐音向众神挑战；}
\verseline{}{妒忌的特里同——若这传说值得相信——抓住他，}
\verseline{}{把他沉进礁石间翻涌的泡沫。\footnote{米塞努斯先后侍奉赫克托耳与埃涅阿斯，以海螺激励战斗。特里同因受到挑战而把他溺死，但叙述特意加上“若这传说可信”的保留语。意大利的米塞努斯岬由他得名，他的死亡也为这一地名提供了神话解释。}}
\verseline{175}{于是众人围着他放声哀号，}
\verseline{}{虔敬的埃涅阿斯尤其悲恸。继而，他们赶办西比尔的命令，}
\verseline{}{毫不拖延，流着泪争相搬来树木，}
\verseline{}{堆起高耸入天的葬礼柴堆。}
\verseline{}{他们走进古老的森林，走入野兽幽深的栖处；}
\verseline{180}{油松倒下，斧头击中的圣栎发出回声，}
\verseline{}{白蜡木梁与易劈的栎木被楔子剖开，}
\verseline{}{他们把巨大的梣木从山上滚下。}
\verseline{}{在这些劳作中，埃涅阿斯也走在最前，}
\verseline{}{一面鼓励同伴，一面备起同样的工具。}
\verseline{185}{他在悲伤的心中反复思量这些事情，}
\verseline{}{望着无边的森林，忽然这样祈求：}
\verseline{}{“但愿那根树上的金枝现在便}
\verseline{}{在这浩瀚的林中向我们露面！毕竟女先知说到你的那些话，}
\verseline{}{唉，米塞努斯啊，竟无一不真。”}
\verseline{190}{他话音刚落，两只鸽子恰好}
\verseline{}{从天空飞到这人面前，}
\verseline{}{停在青翠的地面；伟大的英雄}
\verseline{}{认出母亲的神鸟，欣然祈祷：}
\verseline{}{“请为我引路——若真有路——请在空中指引行程，}
\verseline{195}{领我进入树林；丰饶的枝条正在那里}
\verseline{}{遮蔽肥沃的土地。女神母亲啊，}
\verseline{}{莫在疑难中离弃我。”说完，他停步凝望，}
\verseline{}{察看它们带来什么征兆，又准备向哪里飞去。}
\verseline{}{它们一面觅食，一面振翼前进，始终只飞到}
\verseline{200}{追随者的目光所能守住的距离。}
\verseline{}{等它们来到气味浓重的阿韦尔努斯峡口，}
\verseline{}{便迅速升起，滑过澄澈的空气，}
\verseline{}{在那棵生着两重枝叶的树上，栖在中意之处；金色异光}
\verseline{}{从那里的枝条间闪耀。}
\verseline{205}{正如冬日严寒的森林里，槲寄生常以}
\verseline{}{新叶青翠——它并非由所依附的树木播种——}
\verseline{}{又以淡黄的果实缠绕圆润的树干；}
\verseline{}{金叶在浓荫的圣栎上也是这般景象，}
\verseline{}{薄薄的金片也这样在微风中轻响。}
\verseline{210}{埃涅阿斯立刻抓住金枝，急切地折断}
\verseline{}{那略带迟疑的枝条，带回西比尔的洞府。\footnote{两只鸽子是维纳斯的鸟，为她的儿子引路。槲寄生在冬日仍有绿叶和淡黄果实，又生在别种树上，与圣栎上金枝异色的景象相映。西比尔曾说金枝会自愿、轻易地随受命者而去（146—148行），此处却写它“迟疑”，与埃涅阿斯的急切相对；古代已有以英雄的急切或黄金的柔韧来解释这一“迟疑”的读法，见塞尔维乌斯（6.211）。}}
\verseline{}{与此同时，特乌克罗斯人仍在海岸为米塞努斯}
\verseline{}{哭泣，向不再有感的亡灰献上最后的礼物。}
\verseline{}{他们先用含脂的松木与劈开的栎木}
\verseline{215}{搭起巨大的柴堆，把黑叶编进四壁，}
\verseline{}{又在前面竖起丧葬的柏树，把死者}
\verseline{}{闪亮的武器置于堆顶，作为装饰。}
\verseline{}{一些人备好热水和在火上沸腾的铜釜，}
\verseline{}{为已经冰冷的身体洗浴、涂油。}
\verseline{220}{哀声四起。随后他们把众人哭悼过的遗体放上卧榻，}
\verseline{}{又将紫色衣袍——死者熟悉的遮覆——抛在上面。}
\verseline{}{一些人走到巨大的灵床下，}
\verseline{}{担起这悲伤的职分；他们依照祖先的习俗，}
\verseline{}{背转身去，托着火炬点燃柴堆。堆上的乳香祭品、食物和}
\verseline{225}{盛满橄榄油的调酒钵一同在火中焚烧。}
\verseline{}{灰堆坍落，火焰止息以后，}
\verseline{}{他们以酒洗过遗骨与吸满酒液的余烬；}
\verseline{}{科律奈乌斯把拣起的白骨收入铜瓮。}
\verseline{}{又亲自捧着净水，三次绕过同伴，}
\verseline{230}{以吉祥橄榄枝洒下轻柔的水滴，}
\verseline{}{净化众人，说出最后的告别。}
\verseline{}{虔敬的埃涅阿斯为他垒起一座巨冢，}
\verseline{}{又把死者的武器、船桨和号角安置其上，}
\verseline{}{就在一座临空的山下；此山如今因他}
\verseline{235}{得名“米塞努斯”，世世保有这不朽的名字。\footnote{入冥前先处理意外死亡、尚未安葬的同伴，这一安排近于《奥德赛》11—12 中厄尔佩诺耳请求安葬、随后得到葬礼；米塞努斯墓上的船桨又加强了这一联系。葬礼包含搭建柴堆、洗涤与涂油、覆以紫衣、背身点火、献上酒食、以酒洗灰、收敛遗骨、绕行洒水净化和最后告别。}}
\verseline{}{这些事做完，他迅速执行西比尔的吩咐。}
\verseline{}{那里有一座幽深洞窟，张着广阔可怖的巨口，}
\verseline{}{岩石嶙峋，又有黑湖与森林的阴暗守卫；}
\verseline{}{没有一只飞鸟能在上方安然展翼，}
\verseline{240}{因为从那黑色的喉口涌出的气息}
\verseline{}{直上穹苍，正是这样沉重。}
\verseline{}{因此希腊人称这地方为阿俄耳农——“无鸟之地”。\footnote{“阿俄耳农”来自希腊语“无鸟”之义。古人以洞口喷出的有毒气息解释飞鸟不能越过阿韦尔努斯；Conte 2019 将第242 行以方括号排印，该版校勘记指出若干主要传本缺此行。}}
\verseline{}{女祭司先让四头背脊乌黑的公牛立在那里，}
\verseline{}{把酒倾在它们额上；继而从双角之间}
\verseline{245}{拔取最上端的鬃毛，投入圣火，}
\verseline{}{作为最初的祭礼，同时高声呼唤赫卡忒——}
\verseline{}{她在天上与厄瑞玻斯都掌有权能。}
\verseline{}{其他人把刀置于牲畜喉下，以祭盘承接温热的鲜血。}
\verseline{}{埃涅阿斯亲自挥剑，杀死一只黑羊羔，}
\verseline{250}{献给复仇女神之母和她伟大的姐妹；}
\verseline{}{又为你，普洛塞庇娜，献上一头不孕的母牛。}
\verseline{}{然后他为斯提克斯的君王设下夜祭的祭坛，}
\verseline{}{把公牛的全部血肉放上火焰，}
\verseline{}{又在燃烧的脏器上浇下丰厚的油。\footnote{西比尔先向统辖天上与厄瑞玻斯的赫卡忒呼告。埃涅阿斯又以黑羊献给复仇女神之母“夜”及其伟大的姐妹“大地”，以不孕母牛献给普洛塞庇娜，并为斯提克斯之王狄斯设夜间祭坛。祭牲的颜色、性别和生育状态都在诗中逐一注明。}}
\verseline{255}{看哪，就在朝阳初升、曙光初现之际，}
\verseline{}{脚下的大地开始轰鸣，林木覆盖的山脊开始震动；}
\verseline{}{女神临近，犬吠仿佛穿过阴影传来。}
\verseline{}{“远离这里，远离这里，未受祭仪的人！”}
\verseline{}{女先知高声喊道，“全都退出这片树林；}
\verseline{260}{而你，踏上道路，从鞘中拔出利剑：}
\verseline{}{现在需要勇气，埃涅阿斯，需要一颗坚强的心。”}
\verseline{}{她只说了这些，便狂烈地冲入敞开的洞口；}
\verseline{}{他紧跟向导的步伐，脚步同样毫无畏惧。}
\verseline{}{掌管亡灵的诸神、沉默的幽影啊，}
\verseline{265}{还有混沌、佛勒革同与大片沉入夜静的地方，}
\verseline{}{愿神律准许我说出听闻的一切；愿你们的神威}
\verseline{}{容我揭开深埋在大地与昏暗中的事物！\footnote{本卷第 264—267 行忽然转为叙述者发言。诗人先向统治亡魂的诸神和幽暗本身祈求许可，才敢揭露地下深处的事物；这番祈请把接下来的描写标为通常对生者秘而不宣的知识。}}
\verseline{}{他们在孤独的夜色中幽暗前行，穿过阴影，}
\verseline{}{穿过狄斯空寂的宫室与虚无的国度；}
\verseline{270}{如同月影不明、只吝啬地投下微光时，}
\verseline{}{人在森林里行路；那时朱庇特以阴影遮住天空，}
\verseline{}{黑夜从万物身上夺去了颜色。}
\verseline{}{就在门厅之前，在奥尔库斯最初的咽喉处，}
\verseline{}{哀恸与复仇的忧虑安下卧榻；}
\verseline{275}{苍白的疾病和凄凉的老年居住在那里，}
\verseline{}{还有恐惧、诱人作恶的饥饿与丑陋的贫困——}
\verseline{}{一群望之可怖的形体；还有死亡与劳苦；}
\verseline{}{继而是死亡同族的沉睡，心灵邪恶的}
\verseline{}{欢悦，和守在对面门槛上、带来死亡的战争；}
\verseline{280}{复仇女神铁铸的寝室，和失去理智的纷争——}
\verseline{}{蛇发中束着染血的圣带。}
\verseline{}{中央有一株幽暗的巨榆，展开枝叶与衰老的臂膀；}
\verseline{}{据说虚妄的梦成群把它当作居所，}
\verseline{}{紧紧依附在每一片叶下。}
\verseline{285}{此外还有许多形貌各异的兽类怪物：}
\verseline{}{半人马栖在门边，双形的斯库拉诸影，}
\verseline{}{百臂的布里阿瑞俄斯，以及勒耳那的怪兽，}
\verseline{}{发出可怕尖啸；还有武装着火焰的喀迈拉，}
\verseline{}{戈耳工、哈耳庇厄与三身的怪影。\footnote{哀痛、悔恨、疾病、衰老、恐惧、饥饿、贫困、死亡、劳苦、睡眠与战争等抽象力量被安置在冥府门前，继而出现半人马、斯库拉、布里阿瑞俄斯、勒耳那水蛇、喀迈拉、戈耳工与哈耳庇厄等怪物。西比尔及时提醒埃涅阿斯：它们只是没有实体的幽影，刀剑伤不到它们。}}
\verseline{290}{埃涅阿斯骤然惊惧，立刻拔剑，}
\verseline{}{以出鞘的锋刃迎向来者；}
\verseline{}{若非博识的同伴提醒他，那些纤薄无体的幽魂}
\verseline{}{只在空洞的形影里飘飞，}
\verseline{}{他便会扑上前去，以刀刃徒然劈砍幽影。}
\verseline{295}{从这里有一条路通向塔耳塔洛斯的阿刻戎河水。}
\verseline{}{浑浊的漩涡夹着泥浆，在广阔的深渊里}
\verseline{}{翻腾，把所有泥沙喷吐进科库托斯河。}
\verseline{}{一个可怖的摆渡者守着这些水域与河流：}
\verseline{}{卡戎遍体骇人的污秽，下巴垂着大片}
\verseline{300}{蓬乱的白须；双眼凝立，燃烧如火，}
\verseline{}{肮脏的斗篷在肩头打结，向下悬垂。}
\verseline{}{他亲自用长篙驱船，也照料船帆，}
\verseline{}{用一叶锈色小舟运送亡者的躯体；}
\verseline{}{他已经年老，神的老迈之身却仍苍劲而健旺。\footnote{阿刻戎、科库托斯与斯提克斯在这里组成彼此汇流的冥河水系，原诗并不建立一张完全一致的地下地图。卡戎外表肮脏衰老，诗中却说他的老年仍有神的强健。}}
\verseline{305}{所有亡者的群集都涌到这里，铺满河岸：}
\verseline{}{有母亲、男人，有已经走完生命的\footnote{第 306—308 行几乎逐字重用《农事诗》4.475—477 俄耳甫斯下冥府时的亡魂名列；河岸亡魂聚集的外层场景又承接《奥德赛》第十一卷的冥府群魂。}}
\verseline{}{豪杰之躯，也有少年和未婚的少女，}
\verseline{}{还有在父母眼前被放上柴堆的青年；}
\verseline{}{如初寒降临时，秋叶在森林里纷纷}
\verseline{310}{坠落；又如无数飞鸟从深海聚向陆地，}
\verseline{}{当寒冷的季节驱使它们越过海面，}
\verseline{}{把它们送往阳光温暖的土地。}
\verseline{}{他们站着哀求，要最先渡过这段航程，}
\verseline{}{向彼岸伸出双手，渴望着那片河岸。}
\verseline{315}{阴郁的船夫却一时接纳这些，一时接纳那些，}
\verseline{}{又把另一些远远赶开，不许靠近沙岸。}
\verseline{}{埃涅阿斯惊异于眼前的骚动，心中震动，}
\verseline{}{便问：“少女啊，请告诉我，为何众魂聚集在河边？}
\verseline{}{这些灵魂在求什么？又依据怎样的区分，}
\verseline{320}{一些离开河岸，另一些以桨划过铅灰的浅滩？”}
\verseline{}{年迈的女祭司简短地回答他：}
\verseline{}{“安喀塞斯之子、众神血脉最真确的后裔啊，}
\verseline{}{你看见的是科库托斯的深潭与斯提克斯沼泽；}
\verseline{}{众神也敬畏它的神威，不敢背弃在它面前立下的誓言。}
\verseline{325}{你眼前这整个人群，都是无所凭依、未得安葬的亡者；}
\verseline{}{那个摆渡者是卡戎；被河水载走的，都已入土。}
\verseline{}{在遗骨安息于自己的居所之前，谁也不得}
\verseline{}{渡过这可怖的河岸与嘶哑的流水。}
\verseline{}{他们要在这些岸边飘飞游荡整整百年；}
\verseline{330}{到那时才获准重返向往的水潭。”\footnote{已安葬者获准渡河，未葬者须在岸边游荡一百年，此后才获准。帕利努鲁斯请求埃涅阿斯带他渡河，西比尔不许他违背冥界的规定。}}
\verseline{}{安喀塞斯之子停下来，止住脚步，}
\verseline{}{心中思绪万千，怜悯他们不公的命数。}
\verseline{}{他在那里看见悲哀而未得死亡荣礼的}
\verseline{}{琉卡斯庇斯和吕喀亚舰队的首领俄戎忒斯；}
\verseline{335}{他们一同离开特洛伊，越过多风的海面时，}
\verseline{}{南风将船与人卷进海水，尽数淹没。\footnote{琉卡斯庇斯与吕喀亚船长俄戎忒斯都是第一卷风暴中的遇难者。第一卷只简短写到这些死者；到第六卷的冥河岸边，他们作为无葬亡魂重新出现。}}
\verseline{}{看哪，舵手帕利努鲁斯正向他们走来；}
\verseline{}{不久前在利比亚海程中，他还守望星辰，}
\verseline{}{却从船尾坠落，抛在大海中央。}
\verseline{340}{重重阴影遮住他悲哀的形貌，埃涅阿斯勉强认出，}
\verseline{}{便先开口：“帕利努鲁斯，是哪位神明}
\verseline{}{把你从我们身边夺走，沉入大海中央？}
\verseline{}{说吧。阿波罗从前从未显得虚假，}
\verseline{}{唯独这个答复欺骗了我的心：}
\verseline{345}{他曾预言你会安然渡海，}
\verseline{}{抵达奥索尼亚的疆界。看哪，这就是诺言的信义吗？”}
\verseline{}{他答道：“安喀塞斯之子、我的首领，}
\verseline{}{福玻斯的神谕并没有骗你，也没有神把我沉入海中。}
\verseline{}{只是我奉命紧守、用来指引航程的舵柄}
\verseline{350}{偶然被一股巨力扯断；我依旧攀住它，}
\verseline{}{坠落时把它一同拖进海里。我指着狂暴的大海起誓：}
\verseline{}{我对自身安危并无这般恐惧；我更怕}
\verseline{}{你的船失去舵具，又被震落了舵手，}
\verseline{}{会在如此高涌的波浪间支持不住。}
\verseline{355}{凶猛的南风把我载过无边海面，}
\verseline{}{整整漂流了三个风暴之夜；直到第四天的曙光，}
\verseline{}{我才从高高的浪峰上望见意大利。}
\verseline{}{我一点点游向陆地，已经抓住了安全之处；}
\verseline{}{若非那残酷的族人见我湿衣沉重，}
\verseline{360}{双手如钩，正攀住山崖崎岖的尖顶，}
\verseline{}{便以刀剑袭来，无知地把我当作可以掠取的财物。}
\verseline{}{如今海浪占有我，海岸的风仍翻动我的身体。}
\verseline{}{我凭天空悦人的光与空气，}
\verseline{}{凭你的父亲，也凭正在成长的尤卢斯所寄托的希望求你：}
\verseline{365}{无敌者啊，救我脱离这些灾难；或者给我覆上}
\verseline{}{一抔泥土——你能做到——并前往维利亚港；}
\verseline{}{或者，若真有一条路，若你的女神母亲}
\verseline{}{已经向你显示——我相信，若无众神神意，}
\verseline{}{你绝不会准备渡过如此浩大的河流与斯提克斯沼泽——}
\verseline{370}{就向不幸的我伸出右手，带我一同越过河水，}
\verseline{}{让我至少能在死后安息于宁静的居所。”}
\verseline{}{他说完这些，女先知便回答：}
\verseline{}{“帕利努鲁斯，如此凶险的欲望从何而来？}
\verseline{}{你未得安葬，竟想看见斯提克斯的水和}
\verseline{375}{复仇女神森严的河流，未经命令便走近河岸？}
\verseline{}{别再指望祈祷能扭转众神的命定。}
\verseline{}{但请记住我这番话，把它当作严酷遭际的安慰。}
\verseline{}{邻近的居民将在远近各城受到天降异兆驱使，}
\verseline{}{为你的遗骨赎罪，}
\verseline{380}{为你筑起坟墓，向墓冢献上年年重行的祭礼；}
\verseline{}{那地方也将永远保有帕利努鲁斯之名。”}
\verseline{}{这番话驱散他的忧虑，暂时把悲痛}
\verseline{}{赶出凄苦的心；他欣喜那片土地将以自己得名。\footnote{帕利努鲁斯纠正了第五卷末埃涅阿斯的误解：他连同舵柄坠海后仍在浪中游了三夜，第四天抵达意大利，随即遭当地人杀害并被弃尸。西比尔受冥河法则约束，却预告当地人会为他建墓，并以他的名字命名海岬；迟来的地名与祭祀补偿了眼前无法完成的安葬。帕利努鲁斯在第348行说“也没有神把我沉入海中”；Conte 2019 指出，这与第五卷睡神将他推入海中的叙述并不一致。}}
\verseline{}{于是他们走完已经开始的路，接近河流。}
\verseline{385}{船夫早已从斯提克斯水面看见他们}
\verseline{}{穿过寂静的树林，正把脚步转向河岸，}
\verseline{}{便抢先开口，主动斥责：}
\verseline{}{“无论你是谁，竟携带武器走向我的河流，}
\verseline{}{就站在那里说清来意，停住脚步。}
\verseline{390}{这里是幽影、睡眠与昏沉夜色的居所；}
\verseline{}{以斯提克斯之舟运送活人身体，有违神律。}
\verseline{}{从前阿尔喀得斯来到湖边，我接纳了他，}
\verseline{}{并没有因此欢喜；接纳忒修斯和皮里托俄斯也是一样，}
\verseline{}{尽管他们生自众神，力量无人能胜。}
\verseline{395}{前一个亲手捕捉塔耳塔洛斯的守门者，要把它锁住，}
\verseline{}{又从地下君王本人的宝座前拖走那战栗的怪兽；}
\verseline{}{后两人则企图从狄斯的寝宫带走女主人。”}
\verseline{}{安弗律索斯的女先知简短回答：}
\verseline{}{“这里没有那样的诡计；别再动怒；}
\verseline{400}{这些兵器并不带来暴力。那巨大的守门者尽可}
\verseline{}{在洞里永远狂吠，惊吓没有鲜血的幽影，}
\verseline{}{贞洁的普洛塞庇娜也尽可守住叔父的门槛。}
\verseline{}{特洛伊人埃涅阿斯，以虔敬与武功著称，}
\verseline{}{正下降到厄瑞玻斯最深的幽影中寻找父亲。}
\verseline{405}{若这样深的虔敬仍不能打动你，}
\verseline{}{至少认出这根金枝吧。”她揭开衣袍，露出藏在其中的枝条。}
\verseline{}{他因愤怒胀起的心随即平复。}
\verseline{}{他们没有再多说。卡戎惊望那可敬的礼物，}
\verseline{}{惊望许久未曾看见的命定金枝，}
\verseline{410}{便掉转深蓝的小舟，靠近岸边。}
\verseline{}{随后把坐在长凳上的其他灵魂}
\verseline{}{全都赶下船，腾空舱板；同时让}
\verseline{}{高大的埃涅阿斯登进舱中。缝合的小舟在重压下}
\verseline{}{呻吟，裂缝吞进大量沼泽的水。}
\verseline{415}{他们终于安然渡过河流；卡戎把女先知与英雄}
\verseline{}{放在丑陋的淤泥和灰绿色的莎草间。\footnote{卡戎曾载过赫尔库勒斯、忒修斯和皮里托俄斯；这些活人都曾强行进入冥界，因此他敌视生者。“安弗律索斯的先知”以阿波罗曾在色萨利安弗律索斯河畔牧放牛群的故事称西比尔为阿波罗的女先知。金枝使卡戎息怒；埃涅阿斯的活人重量又让缝合的船体呻吟、进水，具体显出他与轻盈亡魂的差别。}}
\verseline{}{巨大的刻耳柏洛斯横卧在对面的洞里，}
\verseline{}{以三个喉咙的狂吠震响这片国度。}
\verseline{}{女先知看见它颈上的蛇已经根根竖起，}
\verseline{420}{便抛去一块浸着蜂蜜、拌有安眠药草的饼。}
\verseline{}{它在狂烈的饥饿中张开三个喉口，}
\verseline{}{攫住投来的食物；巨大的脊背随即松弛，}
\verseline{}{瘫在地上，庞大身躯铺满整个洞穴。\footnote{守门犬刻耳柏洛斯没有被埃涅阿斯以武力制服。西比尔把掺有催眠药的蜂蜜糕抛给它，它吞下诱饵后倒地沉睡。蜂蜜糕也见于古代墓葬与冥界仪式，诗中以祭品和药物相结合的方式让英雄安全通过。}}
\verseline{}{守卫沉睡以后，埃涅阿斯迅速穿过入口，}
\verseline{425}{迅速离开那条一去不能回返的水流岸边。}
\verseline{}{立刻传来人声与响亮的啼哭；}
\verseline{}{那是婴儿的灵魂，正在最初的门槛上哭泣。}
\verseline{}{他们无缘尝到生命的甜美，才离开乳房，}
\verseline{}{黑暗的日子便夺走他们，沉进苦涩的夭亡。}
\verseline{430}{他们近旁，是受诬告而被判死的人。}
\verseline{}{这些居所也并非未经抽签、未有法官便分派：}
\verseline{}{审判者弥诺斯摇动签瓮；他召集沉默者的}
\verseline{}{会议，查问他们的一生与所受控告。}
\verseline{}{随后的一带由悲哀者占据；他们本无罪过，}
\verseline{435}{却亲手为自己造成死亡，厌弃日光，}
\verseline{}{抛舍了灵魂。如今，他们何等愿意在高天之下}
\verseline{}{忍受贫困，忍受艰苦的劳作！}
\verseline{}{神律却挡住他们；那片可憎流水的凄凉沼泽}
\verseline{}{把他们缚住，横隔其间的斯提克斯以九重河道围困他们。}
\verseline{440}{离这里不远，人们可以看见原野向四面铺开；}
\verseline{}{他们称它为“哀悼之野”。}
\verseline{}{这里，那些被严酷的爱残忍地耗蚀至死的人，}
\verseline{}{藏身在隐秘的小径；四周的桃金娘}
\verseline{}{树林遮蔽着他们；忧苦即使在死后也不肯离开。}
\verseline{445}{他在这里看见斐德拉、普洛克里斯和悲哀的厄里斐勒——}
\verseline{}{她正指着残酷儿子留下的伤口——}
\verseline{}{又看见欧阿德涅和帕西法厄；拉俄达弥亚}
\verseline{}{伴她们同行，还有从前是青年、如今是女子的凯纽斯，}
\verseline{}{命运又使她回归旧日的形体。\footnote{诗中按死亡方式和遭际把亡魂分置在冥界不同地方：婴儿、蒙冤受刑者、自尽者和为爱情耗损者各处一域。悲恋者中列有斐德拉、普洛克里斯、厄里斐勒、欧阿德涅、帕西法厄、拉俄达弥亚和凯纽斯；最后一人曾由女子变为男子，死后又恢复最初的女性形体。}}
\verseline{450}{在她们中间，腓尼基人狄多伤口犹新，}
\verseline{}{正在广大的森林中游荡。特洛伊英雄一走到近旁，}
\verseline{}{便隔着阴影认出她朦胧的身影，}
\verseline{}{正如有人在月初透过云层，}
\verseline{}{看见或以为自己看见一弯初升的新月；}
\verseline{455}{他垂下泪水，声音里带着温柔的爱：}
\verseline{}{“不幸的狄多，传到我这里的消息竟是真的——}
\verseline{}{你已死去，以刀刃追随了最后的结局？}
\verseline{}{唉，我是你死亡的原因吗？我凭群星起誓，}
\verseline{}{凭天上的众神，也凭地下深处若还存在的信义，}
\verseline{460}{王后啊，我是身不由己才离开你的海岸。}
\verseline{}{可是众神的命令如今逼我穿过}
\verseline{}{这些幽影、荒芜霉朽的地方与深不见底的黑夜，}
\verseline{}{也曾以自己的权威驱我离去；我不曾料到}
\verseline{}{我的离去会给你带来这样深的痛苦。}
\verseline{465}{停住脚步，别从我的眼前抽身。}
\verseline{}{你在逃避谁？这是命运让我对你说的最后一番话。”}
\verseline{}{埃涅阿斯用这些话安抚她燃烧的心；她仍怒目相视，}
\verseline{}{他自己却泪流不止。}
\verseline{}{她背过身去，双眼始终钉在地上，}
\verseline{470}{听他开始说话，面容没有丝毫移动，}
\verseline{}{仿佛她只是坚硬的燧石，或矗立的马耳珀索斯岩壁。}
\verseline{}{她终于猛然抽身，背离他而去，}
\verseline{}{进入阴影密布的树林；她从前的丈夫}
\verseline{}{叙凯乌斯在那里回应她的忧苦，以同等的爱爱她。}
\verseline{475}{埃涅阿斯同样被这不公的遭际震动，}
\verseline{}{远远含泪目送，怜悯她离去。\footnote{狄多一幕重新打开第四卷尚未弥合的伤口。埃涅阿斯诉说哀痛、解释自己奉神命离开，并说未曾料到离去会给她带来如此深的痛苦；狄多既不回答，也不回望，最后退回前夫叙凯乌斯身边。这里还有两层醒目的互文：第460行invitus, regina, tuo de litore cessi几乎逐词改写卡图卢斯66.39 invita, o regina, tuo de vertice cessi；卡图卢斯第66首又是卡利马科斯《贝勒尼刻的发卷》的拉丁改写。随后狄多拒不答话、转身离去，重演《奥德赛》11.543—567中埃阿斯拒绝奥德修斯和解的场面。}}
\verseline{}{随后，他吃力地走上获准的道路。他们已经来到}
\verseline{}{最远处那片幽僻的原野，战争中成名者聚集其中。}
\verseline{}{提丢斯在这里迎面而来；这里还有以武艺闻名的}
\verseline{480}{帕耳忒诺派俄斯，以及苍白的阿德剌斯托斯之影；}
\verseline{}{这里还有在人世深受哀悼、从战争中倒下的}
\verseline{}{达耳达诺斯人。埃涅阿斯望见他们列成长队，}
\verseline{}{便叹息起来：格劳库斯、墨冬、忒耳西洛科斯，}
\verseline{}{安忒诺耳的三个儿子，还有献身刻瑞斯祭仪的波吕玻忒斯，}
\verseline{485}{以及伊达乌斯——在这里仍执掌战车，仍手握武器。}
\verseline{}{亡魂成群站在他左右；}
\verseline{}{见他一面并不足够，他们愿长久停留，}
\verseline{}{同他并肩而行，得知他为何来到这里。}
\verseline{}{达那俄斯人的首领和阿伽门农的战阵}
\verseline{490}{一看见他本人，看见武器在阴影中闪耀，}
\verseline{}{便因巨大的恐惧惊慌；一些转身逃走，}
\verseline{}{仿佛从前逃向船只；另一些刚想扬声呼喊，}
\verseline{}{微弱的声音便在张开的嘴边消散。\footnote{特洛伊阵亡者围拢埃涅阿斯，希腊亡魂则惊惧退避，有的转身逃跑，有的发出微弱呼喊；他们死后仍保留着战场上的亲疏与恐惧。}}
\verseline{}{就在这里，他看见普里阿摩斯之子得伊福玻斯：}
\verseline{495}{全身被割裂，面容惨遭毁损，}
\verseline{}{脸、双手与太阳穴都被摧残，双耳遭割去，}
\verseline{}{鼻子也被羞辱性的伤口削断。}
\verseline{}{他畏缩着掩藏骇人的刑痕，埃涅阿斯几乎未能认出，}
\verseline{}{终于主动以熟悉的声音向他招呼：}
\verseline{500}{“善战的得伊福玻斯，出自特乌克罗斯崇高的血脉，}
\verseline{}{是谁竟想从你身上索取这样残酷的刑罚？}
\verseline{}{谁竟有权如此待你？最后一夜的传闻}
\verseline{}{告诉我，你杀死无数佩拉斯吉人，精疲力尽，}
\verseline{}{终于倒在尸身混杂的高堆上。}
\verseline{505}{那时我亲自在罗埃特翁海岸为你筑起空冢，}
\verseline{}{又高声三次呼唤你的亡魂。}
\verseline{}{你的名字与武器仍标明那地方；可是朋友啊，我未能}
\verseline{}{找到你，在离开时把你安置在祖国的土地。”}
\verseline{}{普里阿摩斯之子回答：“朋友啊，你没有遗漏什么；}
\verseline{510}{你为得伊福玻斯、为死者幽影尽到了一切葬礼义务。\footnote{得伊福玻斯遭到残酷肢解，面容几乎不能辨认。埃涅阿斯曾在罗埃特翁海岸为他建空墓，三次高声呼唤亡魂，却未能找到遗体；这类墓可以在尸身失落时承接祭礼。得伊福玻斯明确回答，朋友已经尽到一切责任。}}
\verseline{}{是我的命运和那个拉科尼亚女子致命的罪行}
\verseline{}{把我沉进这些灾祸；她留下了眼前的纪念。}
\verseline{}{我们怎样在虚假的欢乐中度过最后一夜，}
\verseline{}{你是知道的；也不得不记得太过清楚。}
\verseline{515}{命定的木马一跃登上帕加马高处，}
\verseline{}{沉重的腹中载来全副武装的步卒；\footnote{马克罗比乌斯（6.2.25）把这两行与恩尼乌斯失传悲剧《亚历山大》残篇并列：恩尼乌斯同样写木马越过帕加马城墙，腹中载着武装士兵。这也是保存该悲剧残句的古代见证。}}
\verseline{}{她假作歌舞，领着发出祭呼、行酒神祭仪的}
\verseline{}{佛里吉亚女子绕行；自己站在正中，举着}
\verseline{}{一支巨大的火炬，从城堡最高处召唤达那俄斯人。}
\verseline{520}{那时我已被忧虑耗尽，又受睡眠重压，}
\verseline{}{不幸的寝室把我收容；我一躺下，}
\verseline{}{甜美深沉的安眠便压住我，极似宁静的死亡。}
\verseline{}{与此同时，我那卓越的妻子移走屋内所有武器，}
\verseline{}{又从我头下抽走那柄可信赖的剑；}
\verseline{525}{她把墨涅拉俄斯召进屋内，为他打开门槛，}
\verseline{}{显然指望这会成为献给爱人的厚礼，}
\verseline{}{也以为这样便能抹去以往罪行留下的恶名。}
\verseline{}{何必再拖延？他们闯入寝室；还有一个同伴加入，}
\verseline{}{那埃俄卢斯的后裔，正是罪行的怂恿者。众神啊，}
\verseline{530}{若我以虔诚的口祈求惩罚，就把同样的灾祸重施于希腊人！\footnote{得伊福玻斯叙述海伦怎样出卖了他：她移走他的兵器，又召墨涅拉俄斯入室。“埃俄卢斯的后裔”是他对乌利塞斯的辱称，联系乌利塞斯被说成西叙福斯或埃俄卢斯后裔的谱系传说。}}
\verseline{}{可是轮到你了；告诉我，是怎样的遭际}
\verseline{}{把活着的你带到这里？是大海的漂泊驱使你前来，}
\verseline{}{还是众神的告诫？又是哪一种命运在磨折你，}
\verseline{}{让你走到这些无日光的凄苦宫室、混沌之地？”}
\verseline{535}{他们这样轮番交谈时，奥罗拉的玫瑰色四马车}
\verseline{}{已经在以太的航程中越过中天；}
\verseline{}{也许他们会把命定的全部时间都耗在这些话上，}
\verseline{}{但同行者出声提醒，西比尔简短说道：}
\verseline{}{“黑夜奔落了，埃涅阿斯；我们却以哭泣消磨时辰。}
\verseline{540}{就在这里，道路分裂成左右两条：}
\verseline{}{右路通向伟大狄斯的城墙之下，}
\verseline{}{我们由此走向厄吕西乌姆；左路则让恶人}
\verseline{}{承受刑罚，把他们送进邪恶的塔耳塔洛斯。”}
\verseline{}{得伊福玻斯答道：“伟大的女祭司，别再动怒；}
\verseline{545}{我会离开，回到亡魂的行列，归还给黑暗。}
\verseline{}{去吧，去吧，我们的荣光；愿你享有更好的命运！”}
\verseline{}{只说了这些，他便随着话音转过脚步。}
\verseline{}{埃涅阿斯忽然回望，在左侧峭壁之下}
\verseline{}{看见宽广的城垣，四周环绕三重高墙；}
\verseline{550}{塔耳塔洛斯的佛勒革同河卷着湍急的}
\verseline{}{烈焰绕城奔流，翻滚轰响的岩石。}
\verseline{}{正面有一道巨门，门柱由整块金刚石筑成，}
\verseline{}{没有人间的力量，连天上的神也不能}
\verseline{}{在战争中把它摧毁；一座铁塔直立高空，}
\verseline{555}{提西福涅坐在那里，血红的长袍束在腰间，}
\verseline{}{不眠不休，日夜守着门厅。}
\verseline{}{从那里可以听见呻吟，残酷的鞭打}
\verseline{}{不断响起；接着是铁器的铿锵与拖动的锁链。}
\verseline{}{埃涅阿斯停住脚步，惊惧地听尽那片喧响。}
\verseline{560}{“少女啊，请说，那里有怎样的罪行？他们受到}
\verseline{}{怎样的刑罚逼迫？何以如此号哭，直达高空？”}
\verseline{}{女先知于是开口：“特乌克罗斯人显赫的首领，}
\verseline{}{清白之人不得踏上那罪恶的门槛；}
\verseline{}{但赫卡忒让我主掌阿韦尔努斯的树林时，}
\verseline{565}{亲自把众神的刑罚教给我，领我走遍一切。}
\verseline{}{克诺索斯的拉达曼图斯统治这极为严酷的国度，}
\verseline{}{一面惩处，一面审问欺诈，迫使亡魂供认：}
\verseline{}{每个人在上界犯下什么罪，曾徒然庆幸}
\verseline{}{能够隐瞒，把应受的赎罚拖到迟来的死亡。}
\verseline{570}{复仇者提西福涅束好鞭子，立刻扑向罪人，}
\verseline{}{以鞭抽打凌辱；左手伸出狰狞的蛇，}
\verseline{}{召来她那群凶猛的姐妹。}
\verseline{}{到这时，禁忌的门才在轰响的铰链上}
\verseline{}{尖啸着开启。你看见怎样的守卫坐在门厅，}
\verseline{575}{怎样的面孔看守门槛吗？}
\verseline{}{一条更凶暴的巨蛇在里面安下居所，}
\verseline{}{张着五十个幽黑的巨口。而塔耳塔洛斯自身}
\verseline{}{向下裂开，伸入阴影；它的深度是从此地}
\verseline{}{仰望以太中奥林匹斯高天的两倍。\footnote{埃涅阿斯没有进入塔耳塔洛斯。他从外部只见到堡垒、河流、城门与提西福涅，听到刑具和受刑者的声音；内部惩罚由西比尔转述。塔耳塔洛斯向下延伸的深度是从大地仰望奥林匹斯高度的两倍，夸张的垂直尺度强调其与人间、天界的隔绝。}}
\verseline{580}{大地最古老的子女，提坦一族，}
\verseline{}{被雷霆击落，在最深处翻滚。}
\verseline{}{我也在那里看见阿洛阿达双子，躯体巨大；}
\verseline{}{他们曾企图用双手撕开广阔的天空，}
\verseline{}{把朱庇特从天神的王国推下。}
\verseline{585}{我又看见萨尔摩纽斯正为自己招来残酷的刑罚，}
\verseline{}{当他模仿朱庇特的火光与奥林匹斯的雷鸣：\footnote{Conte 2019 校勘记建议将 dantem poenas 理解为 merentem poenas（招致刑罚）：萨尔摩纽斯因模仿朱庇特的雷霆而受罚。}}
\verseline{}{他驾着四马战车，挥舞火把，}
\verseline{}{穿过希腊各族，穿过厄利斯城的中央，}
\verseline{}{耀然凯旋，为自己索取众神的礼敬——}
\verseline{590}{这个狂人竟以铜器和带蹄马匹的踩踏}
\verseline{}{模仿云雷与无可模仿的霹雳。}
\verseline{}{可是全能的父神从浓云之间}
\verseline{}{掷下真正的雷霆，其中没有火把，也没有松明冒烟的}
\verseline{}{火光；一阵巨大的旋风把他头朝下击落。}
\verseline{595}{还可以看见提堤俄斯，大地这位万物之母养育的人；}
\verseline{}{他的身体铺过整整九犹格尔的土地，}
\verseline{}{一只巨鹫用弯曲的喙不断啄食}
\verseline{}{他不死的肝脏与为刑罚再生的内脏，}
\verseline{}{翻寻着作食，栖居在深深的胸膛里；}
\verseline{600}{血肉时时复生，永远得不到片刻安息。\footnote{受刑者包括被雷击坠落的提坦、企图袭击天界的阿洛阿达双子、模仿朱庇特雷霆的萨尔摩纽斯，以及身体横陈九犹格尔土地、肝脏不断再生又不断被秃鹫啄食的提堤俄斯。永不停止的器官更新把刑罚变成没有终点的时间。}}
\verseline{}{又何必提拉庇泰人、伊克西翁和皮里托俄斯？\footnote{Conte 2019 在第601行后排出一处缺文标记。该版校勘记推测原有一行点名坦塔罗斯，为第602—607行的刑罚描写指明对象，并认为缺漏较可能出于传抄。}}
\verseline{}{在那人头上，一块黑岩仿佛马上就要滑落，}
\verseline{}{悬在那里，仿佛随时就会坠下；高高的宴榻上，金色床架}
\verseline{}{闪耀，王者般奢华的宴席摆在他面前；}
\verseline{605}{复仇女神中最强大的一个横卧近旁，}
\verseline{}{阻止他伸手触碰桌案，}
\verseline{}{又举着火把站起，厉声发出雷鸣般的呵斥。}
\verseline{}{这里关着生前憎恨兄弟的人，}
\verseline{}{殴打父母的人，欺骗受自己庇护的门客的人，}
\verseline{610}{或独自伏在所得财富上，}
\verseline{}{不肯分给亲属的人——他们是人数最多的一群——}
\verseline{}{也有因通奸被杀的人，追随不义兵戈的人，}
\verseline{}{以及不曾畏惧背弃与主人相握之右手的人；}
\verseline{}{他们被封闭在此，等待刑罚。不要追问}
\verseline{615}{他们各受什么刑，也别问怎样的情状或命运使他们沉入这里。}
\verseline{}{一些人滚动巨石，一些人被缚在轮辐上}
\verseline{}{悬垂不下；不幸的忒修斯坐在那里，}
\verseline{}{而且将永远坐着；最可怜的佛勒癸阿斯则告诫}
\verseline{}{所有亡者，穿过阴影大声作证：}
\verseline{620}{‘受了警告，就学会正义，不要藐视众神。’}
\verseline{}{这个人为黄金出卖祖国，给它强加一个}
\verseline{}{有权势的主人；又为报酬立法，为报酬废法；}
\verseline{}{那个人侵入女儿的婚床，举行神律所禁的婚礼；}
\verseline{}{所有人都敢犯骇人的渎神之罪，而且得逞。}
\verseline{625}{即使我有一百条舌头、一百张嘴，}
\verseline{}{声音由铁铸成，也不可能穷尽罪行的全部形态，}
\verseline{}{历数所有刑罚的名字。”\footnote{惩罚清单由神话罪人转向具有罗马社会色彩的罪行：仇恨兄弟、殴打父母、欺骗受保护的门客、独占财富、通奸、发动不义战争、背叛主人、受贿枉法与乱伦；这些罪行集中在家庭、庇护关系、财富与公共权力。马克罗比乌斯（6.1.39）将第621—622行列为维吉尔取法瓦里乌斯·鲁福斯《论死亡》（De morte）的例子：两处同以Vendidit hic起首，并共享fixit leges pretio atque refixit。瓦里乌斯是维吉尔同时代的诗人和友人。}}
\verseline{}{福玻斯年迈的女祭司说完这些，又道：}
\verseline{}{“来吧，走上道路，把已经承担的任务做完；}
\verseline{630}{我们加快脚步。我已经望见由独眼巨人熔炉}
\verseline{}{锻造的城墙，还有正面拱顶下的门；}
\verseline{}{神谕吩咐我们把这份礼物留在那里。”}
\verseline{}{她说罢，两人一同穿过幽暗的道路，}
\verseline{}{迅速越过中间的距离，走近门扉。}
\verseline{635}{埃涅阿斯抢先走到入口，以新汲的}
\verseline{}{净水洒遍身体，把金枝固定在正对的门槛上。}
\verseline{}{这些事终于完成，女神的祭礼已经履行，}
\verseline{}{他们便来到欢悦之地，来到幸福树林}
\verseline{}{宜人的绿野和享福者安居的处所。}
\verseline{640}{这里更广阔的以太覆盖原野，给它披上}
\verseline{}{绛紫的光；这里的人有自己的太阳、自己的星辰。\footnote{完成净罪后，埃涅阿斯把金枝安置在普洛塞庇娜的门槛，才进入福地。那里有自己的空气、日月与星辰；原诗称其为“福地”和“幸福的树林”。}}
\verseline{}{一些人在草地的角力场上舒展肢体，}
\verseline{}{在游戏中较量，在金黄的沙地上摔跤；}
\verseline{}{一些人以脚踏出舞步，唱起歌来。}
\verseline{645}{那位身着长袍的色雷斯祭司也在这里，}
\verseline{}{以七根琴弦奏出彼此有别的音声，应和节拍；}
\verseline{}{时而用手指弹拨，时而以象牙拨片击弦。}
\verseline{}{这里有特乌克罗斯古老的族裔，最美好的子孙，}
\verseline{}{生于更佳时代、心志高远的英雄：}
\verseline{650}{伊卢斯、阿萨剌科斯和特洛伊的始祖达耳达诺斯。}
\verseline{}{埃涅阿斯远远惊望他们空置的武器与战车。}
\verseline{}{长矛直插在地上，解开缰绳的马匹}
\verseline{}{散布原野，四处吃草。生前对战车}
\verseline{}{与武器的喜爱，以及饲养光亮骏马的热情，}
\verseline{655}{在他们葬入大地以后仍旧跟随。}
\verseline{}{看哪，他又望见一些人左右坐在草地上}
\verseline{}{进食，合着歌队欣然唱起献给阿波罗的颂歌，}
\verseline{}{就在芬芳的月桂林中；广阔的厄里达努斯河}
\verseline{}{从上方发源，穿过树林滚滚流去。}
\verseline{660}{这里有为祖国战斗而负伤的众人，}
\verseline{}{有生前保持纯洁的祭司，}
\verseline{}{有虔敬的诗人，曾说出配得福玻斯的话，}
\verseline{}{有发明技艺、使生命更为丰饶的人，}
\verseline{}{也有人以功业使后人记住他们；}
\verseline{665}{所有人的额头都环绕着雪白的圣带。}
\verseline{}{他们聚集在周围，西比尔便向众人说道：}
\verseline{}{她首先问穆塞乌斯，因为最多的人群把他}
\verseline{}{围在中央，仰望他高出众人的双肩：}
\verseline{}{“幸福的灵魂们，还有你，最出色的诗人，请告诉我，}
\verseline{670}{哪个地域、哪一处地方留住了安喀塞斯？我们正为他}
\verseline{}{来到这里，渡过了厄瑞玻斯浩大的河流。”\footnote{身着长袍的色雷斯祭司指俄耳甫斯，他以七弦里拉琴应和歌舞。穆塞乌斯身材高大，群魂围绕着他，他又为来客指路。厄里达努斯河流过月桂林；英雄们在福地没有固定居所。}}
\verseline{}{这位英雄用几句话作出回答：}
\verseline{}{“我们没有谁拥有固定的居所；我们住在荫翳的树林，}
\verseline{}{住在河岸铺展的草榻和溪水滋润的青翠草地。}
\verseline{675}{但若你们心愿如此，就越过这道山岭；}
\verseline{}{我会把你们引上平缓的路径。”}
\verseline{}{他说完，便走在前面，从高处为他们指点}
\verseline{}{明亮的原野；随后，他们离开山峰之巅。}
\verseline{}{父亲安喀塞斯却在深处一片青翠的谷地里}
\verseline{680}{细看那些被留在此地、即将重返上界光明的灵魂；}
\verseline{}{他全神贯注，追忆并检阅他们，恰好正把}
\verseline{}{自己的全部族裔与亲爱的后孙逐一清点，}
\verseline{}{也清点这些人的命数、福运、品格与功业。}
\verseline{}{他一看见埃涅阿斯迎面穿过草地走来，}
\verseline{685}{立刻欣喜地伸出双手；}
\verseline{}{泪水从面颊涌落，一声呼唤脱口而出：}
\verseline{}{“你终于来了！你那令父亲久久等待的虔敬}
\verseline{}{终于战胜艰难的路途？我终于能端详}
\verseline{}{你的面容，听见熟悉的声音，同你交谈？}
\verseline{690}{我心中正是这样推想，相信事情终会发生，}
\verseline{}{逐日计算时光；我的牵挂没有欺骗我。}
\verseline{}{你越过怎样的陆地、怎样浩瀚的海面，才来到我这里！}
\verseline{}{儿子啊，你被怎样巨大的危难抛来掷去！}
\verseline{}{我曾多么害怕，利比亚的王国会伤害你！”}
\verseline{695}{埃涅阿斯答道：“父亲，是你，是你那悲哀的幻影}
\verseline{}{屡次迎面显现，驱使我来到这道门槛；}
\verseline{}{船队正停在第勒尼安海上。父亲，请伸出右手，}
\verseline{}{让我握住；别从我的怀抱中抽身。”}
\verseline{}{他如此说着，泪如雨下，浸湿了面颊。}
\verseline{700}{三次，他试图用双臂环住父亲的颈项；}
\verseline{}{三次，那徒然被抱住的幻影都逃出双手，}
\verseline{}{宛如轻风，又极似飞掠的梦境。\footnote{本卷第 700—702 行中，埃涅阿斯三次拥抱父亲而三次落空，几乎逐字重复第二卷第 792—794 行拥抱克瑞乌萨的场面。第五卷第 742 行中，安喀塞斯的幻影也曾如烟消散，埃涅阿斯追问是谁阻止父亲进入自己的怀抱；到了冥府，他终于见到父亲，三次拥抱依旧落空。第 702 行把幻影比作轻风与飞掠的梦。}}
\verseline{}{这时，埃涅阿斯望见幽深的谷地中}
\verseline{}{有一片隔绝的树林，灌木在林间发出声响，}
\verseline{705}{还有勒忒河，流过宁静居所的前方。}
\verseline{}{无数族类和民众在河的周围飞动；}
\verseline{}{如同晴朗的夏日，蜜蜂在草地上}
\verseline{}{栖落于各色花朵，围绕雪白的}
\verseline{}{百合纷纷散开；整片原野都因低鸣而喧响。}
\verseline{710}{埃涅阿斯对这景象骤然战栗，一无所知，}
\verseline{}{便问起缘由：远处那些是什么河流，}
\verseline{}{又是怎样的人，竟以如此大群填满河岸。}
\verseline{}{父亲安喀塞斯回答：“命运还欠这些灵魂}
\verseline{}{第二具身体；他们来到勒忒河的水边，}
\verseline{715}{饮下使人无忧的水与漫长的遗忘。}
\verseline{}{这些灵魂，我早已希望为你一一说出，当面指给你看；}
\verseline{}{我早已盼望清点我族未来的后裔，}
\verseline{}{让你能与我一同为寻得意大利而更加欢喜。”\footnote{勒忒河边的灵魂将再次获得肉身；他们在规定的千年期满后饮下遗忘之水，才会失去旧生记忆并渴望重返人世。}}
\verseline{}{“父亲，难道真有灵魂会从这里升向天空，}
\verseline{720}{高高离去，随后又回到迟钝的}
\verseline{}{身体？这些不幸者为何如此凶险地渴望光明？”}
\verseline{}{“我会告诉你，儿子，不让你悬心等待。”}
\verseline{}{安喀塞斯接过话，把一切依次展开：}
\verseline{}{“最初，天空、大地、流动的原野，}
\verseline{725}{发光的月球与提坦血统的星辰，都由内在的}
\verseline{}{灵气养育；心神灌注所有肢体，推动}
\verseline{}{整个宇宙，与这庞大的身体混合。}
\verseline{}{由此生出人类、牲畜与飞鸟的生命，}
\verseline{}{也生出大海在大理石般的水面下负载的怪物。}
\verseline{730}{这些生命的种子具有火的活力，源自天上；}
\verseline{}{只是有害的身体拖慢它们，}
\verseline{}{大地所造的肢体和终将死去的血肉使它们迟钝。}
\verseline{}{因此它们恐惧、欲求、悲哀、欢喜，却看不清}
\verseline{}{天空，因为被封闭在黑暗和盲目的牢狱。}
\verseline{735}{即使到了最后一日，生命已离开他们，}
\verseline{}{不幸者的一切恶质仍未尽去，一切}
\verseline{}{肉身的污损也未从根底消散；许多东西长久}
\verseline{}{凝结在深处，必然以奇异的方式积习成性。}
\verseline{}{所以他们受刑磨炼，为旧日的恶}
\verseline{740}{偿付惩罚：一些灵魂无遮无蔽地悬起，}
\verseline{}{任风吹荡；另一些在广阔的深流下}
\verseline{}{洗去浸染的罪垢，或用火焚净；}
\verseline{}{我们每一个都承受各自的幽冥刑罚；此后}
\verseline{}{才被送往广阔的厄吕西乌姆，只有少数人居于欢悦的原野，}
\verseline{745}{直到漫长的年月走完时光的循环，}
\verseline{}{除去凝结的污点，只留下纯净的}
\verseline{}{以太知觉，以及不杂他物的灵气之火。}
\verseline{}{这些灵魂全都转过一千年的轮回以后，}
\verseline{}{神便把他们大群召到勒忒河边；}
\verseline{750}{让他们忘却一切，重新看见上界的穹苍，}
\verseline{}{并再次开始渴望回到身体之中。”\footnote{安喀塞斯先说世界由一种精神贯通，继而说灵魂受肉身污染、死后经过净化，再经轮回返回肉身。principio、inde、hinc 等连接词以及 seminibus、aurai 等词带有卢克莱修《物性论》的教诲诗语言，但这里讲述的并非伊壁鸠鲁派灵魂观。Quisque suos patimur Manes 是本段著名难句，塞尔维乌斯（6.743）已列不止一种解释；结合前后的净罪内容，可理解为“我们各自承受属于自己的冥界惩罚／净炼”。}}
\verseline{}{安喀塞斯说完，便把儿子和西比尔一同}
\verseline{}{带进集会中央，走入喧响的人群；}
\verseline{}{又登上一座土丘，从那里可以沿长长的队列}
\verseline{755}{细看所有迎面而来者，辨清渐近的面容。}
\verseline{}{“现在听吧：达耳达诺斯的后裔将有怎样的}
\verseline{}{荣光接踵而来，意大利族裔中有哪些子孙在等待，}
\verseline{}{有哪些显赫的灵魂将继承我们的名号；}
\verseline{}{我会一一说明，也把你自己的命运教给你。}
\verseline{760}{你看，那个倚着一支无刃长矛的青年，}
\verseline{}{命数给他最靠近光明的位置；他将第一个升向}
\verseline{}{以太的空气，血脉中混入意大利人的血——}
\verseline{}{西尔维乌斯——一个阿尔巴人的名号——你死后出生的儿子。}
\verseline{}{你年事已高时，妻子拉维尼娅才会}
\verseline{765}{在树林中把他养大；他将是国王，也是众王之父，}
\verseline{}{我们的族裔将由他开始统治阿尔巴隆伽。\footnote{西尔维乌斯是埃涅阿斯与拉维尼娅之子，名字与“森林”相连。postuma 可指晚生或遗腹；longaeuo 可指年老，也有指埃涅阿斯神化后长久生命的古代解释，见盖利乌斯《阿提卡之夜》2.16。阿尔巴王系由他展开，将特洛伊血统连接到罗马；“西尔维乌斯”成为该王系反复沿用的名号。}}
\verseline{}{紧邻他的是普罗卡斯，特洛伊人的荣光，}
\verseline{}{还有卡皮斯、努米托耳，以及使你的名字重现的}
\verseline{}{西尔维乌斯·埃涅阿斯；无论虔敬或武功}
\verseline{770}{同样卓越，只要他有朝一日接掌阿尔巴的王权。}
\verseline{}{看哪，怎样的青年！他们显出何等强大的力量，}
\verseline{}{额头上又戴着公民橡叶投下的荫影！}
\verseline{}{他们将为你建立诺门图姆、加比和菲得奈城，}
\verseline{}{把科拉提亚的堡垒安置在山上，}
\verseline{775}{还将建起波墨提伊、伊努乌斯堡、波拉与科拉。}
\verseline{}{那时这些将成为名字；如今仍是无名的土地。\footnote{普罗卡斯、卡皮斯、努米托耳和西尔维乌斯·埃涅阿斯属于阿尔巴王系。额上饰有栎叶的“公民冠”在罗马授予救过同胞性命者；安喀塞斯又列出诺门图姆、加比、菲得奈、科拉等未来城市，仿佛这些尚未建成的城邑已经可以看见。}}
\verseline{}{玛尔斯之子罗慕路斯也会加入，陪伴他的外祖父；}
\verseline{}{他的母亲伊利娅出自阿萨剌科斯的血脉，}
\verseline{}{将把他养大。你可看见他头顶竖起的双重冠羽}
\verseline{780}{也看见他的父亲已用自己的神圣荣耀标记了他吗？}
\verseline{}{看哪，儿子，在他的吉兆下，声名卓著的罗马}
\verseline{}{将使统治权与大地等广，使气魄与奥林匹斯等高，}
\verseline{}{又以一道城墙围住自己的七座山丘，}
\verseline{}{因英雄子孙繁盛而幸福；正如柏瑞钦托斯的母神}
\verseline{785}{头戴塔冠，乘车经过佛里吉亚的城市，}
\verseline{}{因生育众神而欢喜，把一百个子孙抱在怀里，}
\verseline{}{他们全都居于天上，全都占据高处的世界。\footnote{罗慕路斯是战神玛尔斯与伊利娅之子，日后又以奎里努斯之名受到神化崇拜。 “柏瑞钦托斯的母神”是佛里吉亚大母神库柏勒，头戴城楼冠，乘车巡视诸城；诗以这位头戴城冠、巡视诸城的母神来比拟罗马七丘及其众多后裔。}}
\verseline{}{现在把双眼转向这里，看这个族裔，}
\verseline{}{看你自己的罗马人。这里有恺撒，也有尤卢斯的}
\verseline{790}{全部后裔，都将来到浩大的天轴之下。}
\verseline{}{这个人，正是这个人！你一再听见预言所许的}
\verseline{}{奥古斯都·恺撒，神化者之子；他将在拉提乌姆}
\verseline{}{重建黄金时代——在那片田野，萨图尔努斯从前曾经统治——}
\verseline{}{又将越过加拉曼特斯人与印度人}
\verseline{795}{扩展帝国——那片土地位于星辰之外，}
\verseline{}{在年与太阳的道路之外；擎天的阿特拉斯在那里}
\verseline{}{以肩转动镶嵌着燃烧星辰的天穹。}
\verseline{}{他尚未来临，里海的王国已经因}
\verseline{}{众神的答复而战栗，迈俄提斯的土地也是如此，}
\verseline{800}{七流尼罗河的河口也在惊惶中骚动。\footnote{本卷第 788—800 行颂扬奥古斯都·恺撒： “神化者之子”指他经收养而继承被神化的尤利乌斯·恺撒。预言把奥古斯都的统治说到印度、加拉曼特斯、里海、迈俄提斯湖和尼罗河口，并把他的时代同萨图尔努斯的黄金时代相比；随后却转入内战劝诫和小马塞卢斯的早夭。}}
\verseline{}{就连阿尔喀得斯也未曾走遍如此广大的土地，}
\verseline{}{尽管他射中那铜蹄的母鹿，又让厄律曼托斯的}
\verseline{}{森林恢复安宁，以弓箭使勒耳那战栗；}
\verseline{}{胜利的利柏尔也没有走遍；他以藤蔓为缰驾车，}
\verseline{805}{赶着猛虎，从尼萨高峰而下。}
\verseline{}{难道我们还在犹豫，不肯让勇武在功业中伸展开来，\footnote{virtutem extendere factis 把“勇武”作为伸展的对象，把“行动／功业”作为实现它的方式；这句与下一行合为一个反问。}}
\verseline{}{还是恐惧阻止我们立足奥索尼亚的土地？\footnote{本卷第 801—807 行把奥古斯都的行程置于赫尔库勒斯和利柏尔之上：前者征服铜蹄鹿、厄律曼托斯与勒耳那，后者驾虎车自尼萨山而下。安喀塞斯继而以反问催促特洛伊人，不应再犹疑于让勇武在功业中伸展开来、立足奥索尼亚。}}
\verseline{}{远处那位以橄榄枝为标志、}
\verseline{}{手持圣物的人是谁？我认出了罗马王的}
\verseline{810}{头发与花白的胡须；他将以法律}
\verseline{}{奠定这座初生的城市，从小小的库瑞斯与贫瘠土地}
\verseline{}{被派来掌握浩大的统治权。继他而来的}
\verseline{}{是图卢斯；他将打破祖国的安宁，推动迟缓的}
\verseline{}{男子拿起武器，也推动早已不习惯凯旋的}
\verseline{815}{战阵。安库斯紧随在他身旁，更加自负，}
\verseline{}{即使此刻，也已经太过陶醉于民众的拥戴。}
\verseline{}{你还想看塔克文诸王、复仇者布鲁图斯}
\verseline{}{高傲的灵魂，以及重新夺回的束棒吗？\footnote{手持橄榄枝、准备祭礼的无名国王是努马·庞皮利乌斯；他之后依次出现尚武的图卢斯、好夸耀的安库斯与塔克文诸王。权标束把棍棒与斧头捆在一起，是罗马高级官员命令权的象征；布鲁图斯夺回权标，表示王政终结、共和国建立。}}
\verseline{}{他将最先取得执政官的权力与严酷的斧钺；}
\verseline{820}{身为父亲，却要把发动新战争的儿子}
\verseline{}{召来受刑，以换取光荣的自由。}
\verseline{}{不幸的人啊，无论后世怎样传说这件事，}
\verseline{}{对祖国的爱和对荣誉无边的渴望都将获胜。\footnote{卢基乌斯·尤尼乌斯·布鲁图斯镇压恢复王政的阴谋，并亲自判处参与其中的两个儿子死刑。诗称他“不幸”，又写他对祖国的爱与对荣誉的强烈渴望；赞许因此留有张力。}}
\verseline{}{再看远处的德基乌斯、德鲁苏斯家族，看}
\verseline{825}{执斧严酷的托尔夸图斯和夺回军旗的卡米卢斯。\footnote{德基乌斯父子以把自身献给地下诸神的仪式换取罗马军队胜利；德鲁苏斯家族以战功著称；托尔夸图斯曾因违抗军令处死儿子；卡米卢斯则从高卢人手中夺回军旗。短短几行连续列出自我牺牲、家法、军纪和复国战功。}}
\verseline{}{而那边，你看见两个灵魂以相同的武装闪耀；}
\verseline{}{此刻，只要仍被黑夜压住，他们和睦相处，}
\verseline{}{可是唉，一旦触及生命的光，彼此将唤起}
\verseline{}{何等巨大的战争、怎样的战阵和杀戮！}
\verseline{830}{岳父从阿尔卑斯山的壁垒与摩诺埃库斯的城堡}
\verseline{}{向下而来，女婿则率东方大军在对面列阵。}
\verseline{}{不要，孩子们，不要让心灵熟习这样深重的战争，}
\verseline{}{不要把强大的力量刺向祖国的脏腑；}
\verseline{}{你要率先罢手，你的血脉出自奥林匹斯，}
\verseline{835}{把兵器从手中抛下吧，我的骨血！——\footnote{披着同样铠甲、终将相争的“女婿”和“岳父”是庞培与尤利乌斯·恺撒；恺撒之女尤利娅的婚姻曾连接两人。安喀塞斯恳求他们不要让罗马人的力量反噬自身，并先向作为尤利乌斯家族祖先的埃涅阿斯之“后裔”恺撒发出放下武器的呼告。}}
\verseline{}{那个人征服科林斯以后，将驾胜利的战车}
\verseline{}{登上高高的卡庇托利乌姆，因斩杀阿凯亚人而闻名。}
\verseline{}{另一人将毁灭阿尔戈斯与阿伽门农的迈锡尼，}
\verseline{}{也将击败埃阿科斯的后裔、善战阿喀琉斯的子孙，}
\verseline{840}{为特洛伊的先祖和遭亵渎的弥涅耳瓦神庙复仇。}
\verseline{}{伟大的加图啊，谁会默然略过你，或略过你，科苏斯？}
\verseline{}{谁会略过格拉古家族，或那对西庇阿——}
\verseline{}{战争的两道霹雳、利比亚的毁灭——略过资财微薄却}
\verseline{}{掌握大权的法布里基乌斯，或在犁沟播种的你，塞拉努斯？}
\verseline{845}{法比乌斯家族啊，你们把疲惫的我带往哪里？你就是马克西穆斯，}
\verseline{}{唯有你以拖延为我们挽回国势。\footnote{这组共和国人物包括征服科林斯的穆米乌斯、战胜马其顿王珀耳修斯的埃弥利乌斯·保卢斯、加图、科苏斯、格拉古兄弟、两位“非洲征服者”西庇阿、清贫的法布里基乌斯、田间播种的塞拉努斯与以拖延挽救国家的法比乌斯·马克西穆斯。这些名字排成一串共和国人物谱。这里还叠着两处罗马前代诗歌的原声：“战争的两道霹雳”呼应卢克莱修《物性论》第三卷称西庇阿为 belli fulmen；法比乌斯的“唯有你以拖延为我们挽回国势”则几乎照取恩尼乌斯《编年史》第十二卷 unus homo nobis cunctando restituit rem。这两处诗句的比较，分见马克罗比乌斯（6.1.46、6.1.23）。}}
\verseline{}{其他人会把仿佛呼吸的青铜锻得更加柔和，}
\verseline{}{我确实相信，也会从大理石中引出活人的面容；}
\verseline{}{他们会更善于为讼案陈词，以规尺描绘}
\verseline{850}{天体的运行，说出星辰升起的时刻；}
\verseline{}{罗马人啊，你要记得以统治权治理万民；}
\verseline{}{这将是你的技艺：为和平立下秩序，}
\verseline{}{宽恕服从者，彻底征服骄傲者。”\footnote{安喀塞斯把雕塑、雄辩与天文学让给其他民族，把治理万民、建立和平、宽待归顺者和击败骄傲者规定为罗马的“技艺”。这段关于罗马使命的话紧接此前的内战人物，随后又转入小马塞卢斯的哀歌。}}
\verseline{}{父亲安喀塞斯这样说，又向惊叹的听者补充：}
\verseline{855}{“看哪，马塞卢斯佩着从敌帅身上夺来的兵甲走来，}
\verseline{}{胜利者的身影高出所有人！}
\verseline{}{国家被巨大的动乱震撼时，他将使罗马}
\verseline{}{重新站稳；又将亲自策马击溃布匿人与反叛的高卢人，}
\verseline{}{把第三套夺来的主帅武装悬献给父神奎里努斯。”\footnote{此处首先出现的是老马尔库斯·克劳狄乌斯·马塞卢斯。他亲手杀死高卢首领维里多马鲁斯，将从敌帅身上剥下的兵甲奉献给朱庇特·费雷特里乌斯；这类由罗马统帅亲手杀死敌军统帅后取得的兵甲称为 spolia opima，传统中仅有三例。“父神奎里努斯”即神化的罗慕路斯。史家通常说马塞卢斯把兵甲奉献给朱庇特·费雷特里乌斯；维吉尔此处则写 patri Quirino。}}
\verseline{860}{这时埃涅阿斯开口；他看见与马塞卢斯同行的}
\verseline{}{是一个容貌出众、兵甲闪耀的青年，}
\verseline{}{额头却少有喜色，目光随着面容低垂：}
\verseline{}{“父亲，那个伴随英雄走来的人是谁？}
\verseline{}{是他的儿子，还是伟大族系中的某位后裔？}
\verseline{865}{随从在他周围发出何等喧响！他本人何等庄严！}
\verseline{}{可是幽黑的夜正以悲哀的阴影环绕他的头颅。”}
\verseline{}{父亲安喀塞斯这才开口，泪水已经涌起：}
\verseline{}{“儿子啊，别再追问你族人这场巨大的哀恸；}
\verseline{}{命运只会让大地看见他，不会}
\verseline{870}{容许他停留更久。天神啊，若这份厚赠真能长留人间，}
\verseline{}{罗马的后裔在你们眼中便强盛得过了头！}
\verseline{}{战神伟大城邑近旁的那片原野会发出}
\verseline{}{多少壮士的呻吟；台伯河神啊，当你从新筑的}
\verseline{}{墓冢旁流过，又将看见怎样的葬礼！}
\verseline{875}{伊利乌姆族裔的任何少年都不会使拉丁人的}
\verseline{}{祖先升起如此高的希望，罗慕路斯的土地}
\verseline{}{以后也不会再为任何儿子这样骄傲。}
\verseline{}{唉，他的虔敬！唉，古老的信义与战无不胜的}
\verseline{}{右手！他全副武装时，没有谁能迎面而来}
\verseline{880}{仍不受伤害——无论他徒步走向敌人，}
\verseline{}{还是用马刺刺入喷吐白沫的战马两胁。}
\verseline{}{唉，应受怜悯的少年，但愿你能冲破严酷的命运！}
\verseline{}{你将成为马塞卢斯。献上满手的百合；}
\verseline{}{让我撒下紫色花朵，至少以这些礼物}
\verseline{885}{覆在后裔的灵魂之上，完成这徒然的祭献。”}
\verseline{}{他们就这样漫游整个地域，\footnote{年轻人是小马塞卢斯，奥克塔维娅之子、奥古斯都的外甥与女婿，曾被寄予继承希望，于公元前23年早逝。安喀塞斯先哀叹他的早夭和葬礼，继而称赞虔敬、信义与武勇，最后以献花哀悼。Conte 2019 校勘记将882行的si按愿望式理解（“但愿”），又在883行下注中以少年的天真与死亡的哀痛解释百合、紫色花朵的对照；purpureos flores也有仍指百合的读法。“徒然的祭献”写献花已不能挽回生命。}}
\verseline{}{走过空气广阔的原野，察看一切。}
\verseline{}{安喀塞斯领儿子逐一看过这些以后，}
\verseline{}{又以对未来声名的热爱点燃他的心，}
\verseline{890}{继而告诉英雄此后必须进行的战争，}
\verseline{}{教他认识劳伦图姆的民族和拉提努斯的城，}
\verseline{}{也教他应当怎样躲避、怎样承受每一种劳苦。}
\verseline{}{“睡眠有两扇门；其中一扇据说由角制成，}
\verseline{}{真实的幽影可以从那里轻易出去；}
\verseline{895}{另一扇以洁白的象牙制成，光亮耀眼，}
\verseline{}{亡魂却从这里把虚假的梦境送往天空。”\footnote{Conte 2019 将第 893—896 行解释并排印为安喀塞斯的直接引语；Nauck 与 Ribbeck 曾删疑这四行。随后第 897 行“说完这些”也与这种直接引语的处理相应。}}
\verseline{}{说完这些，安喀塞斯陪着儿子和西比尔，}
\verseline{}{将他们从象牙门送出；\footnote{角门与象牙门承自《奥德赛》19.562—567的梦门传统：由角门通过的梦会应验，由象牙门通过的梦欺骗人。依 Conte 2019 的解释，角门在这里只供真正的亡魂通行；埃涅阿斯和西比尔仍是生者，所以安喀塞斯让他们由象牙门返回人间。}}
\verseline{}{埃涅阿斯径直赶往船只，重见同伴；}
\verseline{900}{继而沿海岸直驶向卡伊埃塔港。\footnote{recto litore 写船只沿海岸径直驶向港口；第七卷开篇随即叙述乳母卡伊埃塔去世、安葬并使这片海岸得名。}}
\verseline{}{锚从船首抛下，船尾沿岸停列。\footnote{Conte 2019 将第 901 行以方括号排印；该行与第三卷第 277 行相同，该版校勘记注明塞尔维乌斯未释，并记录 Ribbeck 的删疑。}}
